% !TeX program = pdflatex
% Compile twice with pdfLaTeX.
\documentclass[11pt,a4paper]{article}
\usepackage[T1]{fontenc}
\usepackage[margin=27mm]{geometry}
\usepackage{lmodern}
\usepackage{amsmath,amssymb,amsthm}
\usepackage{booktabs,longtable,array}
\usepackage[hidelinks,hypertexnames=false]{hyperref}
\urlstyle{same}
\newtheorem{theorem}{Theorem}
\newtheorem{lemma}[theorem]{Lemma}
\setlength{\emergencystretch}{2em}
\allowdisplaybreaks[1]
\title{An Explicit Uniform Gap for the Lonely Runner Problem}
% Add the authors and affiliations before submission.
\author{}
\date{}
\begin{document}
\maketitle
\begin{abstract}
Let \(V\) be a set of \(n\ge2\) distinct positive integer speeds.
We prove that there is a time \(t\) at which
\(\min_{v\in V}\|vt\|>(1+e^{-6888})/(2n)\).
The covering argument of Hu and Jing reduces the problem to a weighted
greatest-common-divisor sum on separated speed sets.
We control the distribution of prime valuations using a one-dimensional
Green kernel, and combine a sieve for rough rational numbers with a dyadic
decomposition to bound the sum. The finite numerical inequalities are
certified by rational computations in a Lean formalization.
\end{abstract}

\section{Introduction}

Write \(\|x\|=\min_{k\in\mathbb Z}|x-k|\) and
\(\mathbb T=\mathbb R/\mathbb Z\). For a finite set \(V\) of positive
integer speeds, define
\[
 \mathrm{ML}(V)=\max_{t\in\mathbb T}\min_{v\in V}\|vt\|.
\]
The lonely runner conjecture predicts that \(\mathrm{ML}(V)\ge 1/(n+1)\)
when \(|V|=n\), corresponding to \(n\) moving runners and one stationary
runner. Hu and Jing~\cite{HuJing} proved that there is an absolute
constant \(c>0\) for which \(\mathrm{ML}(V)>(1+c)/(2n)\) for every
\(n\ge2\). Their argument gives explicit constants, without attempting
optimization. Here we give the value \(c=e^{-6888}\), together with
a self-contained proof of the estimates used below.

\begin{theorem}\label{thm:main}
If \(V\subset\mathbb Z_{>0}\) is finite and \(n=|V|\ge2\), there exists
\(t\in[0,1)\) such that
\[
 |vt-k|>\frac{1+e^{-6888}}{2n}
 \qquad(v\in V,\ k\in\mathbb Z).
\]
\end{theorem}

The reduction uses the covering multiplicities of Hu and Jing and the
valuation-layer approach to large common divisors of Green and
Walker~\cite{GreenWalker}. The rational certificates and their
formalization are described in Appendix~\ref{sec:cert}.

For \(0\le r\le 1/2\) let \(B_v(r)=\{t\in\mathbb T:\|vt\|\le r\}\).
Its normalized measure is \(2r\). Set \(r_0=1/(2n)\).

\section{Circle coverings and spillover}

\subsection{Pairwise overlap}

Following~\cite{HuJing}, we separate overlaps at a common center from
those at distinct centers. Define \(H(u,v)=\max(u,v)/\gcd(u,v)\).
For \(r<1/2\), the arcs forming \(B_v(r)\) have centers \(j/v\) and
radius \(r/v\), and are disjoint. The \(\gcd(u,v)\) common centers
contribute a total overlap \(2r/H(u,v)\). Set
\[
 Q_r(u,v)=|B_u(r)\cap B_v(r)|-\frac{2r}{H(u,v)}.
\]

\begin{lemma}\label{lem:pair}
For positive integers \(u,v\) and \(0\le r<1/2\),
\[
 \max\{0,4r^2-4r/H(u,v)\}\le Q_r(u,v)\le4r^2.
\]
\end{lemma}

\begin{proof}
Write \(u=da\) and \(v=db\), where \(a\le b\) and \(\gcd(a,b)=1\).
The map \(t\mapsto dt\) preserves normalized measure, so it suffices to
use the speeds \(a,b\). The differences \(j/a-k/b\bmod 1\) run once
through the grid of mesh \(1/(ab)\) by coprimality. On rescaling by
\(ab\), the two radii become \(x=br\) and \(h=ar\). Summing over lifts
to the real line, the total overlap multiplied by \(ab\) is
\[
 \sum_{\ell\in\mathbb Z}|[-x,x]\cap[\ell-h,\ell+h]|
 =\int_{-h}^{h}\#\{\ell\in\mathbb Z:-x\le\ell+s\le x\}\,ds .
\]
For \(r>0\), the integrand equals
\(\lfloor x-s\rfloor+\lfloor x+s\rfloor+1\) almost everywhere.
Writing \(\{w\}=w-\lfloor w\rfloor\) and removing the contribution
\(2h\) of the common center gives
\[
 Q_r(u,v)=4r^2-\frac2{ab}\int_{x-h}^{x+h}\{w\}\,dw.
\]
The bounds follow from \(0\le\{w\}<1\) and the nonnegativity of the
overlap at distinct centers. At \(r=0\), all arcs are singletons and the overlap has measure zero.
\end{proof}

\subsection{Dilation of multiplicities}

Let \(N_r=\sum_{v\in V}\mathbf1_{B_v(r)}\), and write
\[
 E_k(r)=\int_{\mathbb T}(N_r-k)_+,\qquad
 g(V)=1-\left|\bigcup_{v\in V} B_v(r_0)\right|,
 \qquad x_+=\max(x,0).
\]
Since \(\int_{\mathbb T}N_{r_0}=1\), we have \(E_1(r_0)=g(V)\).

\begin{lemma}\label{lem:dilation}
For a finite family of arcs dilated about their centers by \(s\ge1\),
the integral of the positive part of the multiplicity minus \(k\)
increases by at least a factor \(s\), provided the dilated arcs have
length less than \(1\). In particular, \(E_k(sr)\ge sE_k(r)\) for every
integer \(k\ge0\) when \(sr<1/2\).
\end{lemma}

\begin{proof}
First consider intervals on the real line. With fixed radii \(\rho_i\),
the length of the union of the intervals centered at \(uc_i\) is a
continuous piecewise affine function of \(u>0\). On an affine piece,
if the extreme endpoints of a component come from intervals \(i,j\),
their extremality gives
\[
 u(c_i-c_j)\le \rho_i-\rho_j,\qquad
 u(c_i-c_j)\le \rho_j-\rho_i.
\]
Thus the slope \(c_j-c_i\) of the component length is nonnegative.
Rescaling the line shows that the union length divided by the dilation
factor is nonincreasing when the radii are dilated at fixed centers.
This applies to every subfamily.

For \(k\ge1\), color the intervals in order of increasing left endpoint
using \(k\) colors. Use an unused color if possible. Otherwise assign
the color whose largest right endpoint so far is smallest. If a new
interval and an earlier interval of the same color contain \(t\),
each other color already has an interval reaching \(t\) from the left.
Consequently the number of colors at a point of multiplicity \(N\)
is \(\min(N,k)\). Any coloring has at most this many colors at a
point. Therefore
\[
 \int(N-k)_+=\sum_i |I_i|-
 \max_{\text{colorings}}\sum_{d=1}^k
       \left|\bigcup_{\operatorname{color}(i)=d}I_i\right|.
\]
The first term grows linearly with dilation, while the second term
divided by the dilation factor is nonincreasing. For \(k=0\), the
integral is the sum of interval lengths.

For the circle, choose one representative of each center in \([0,1)\)
and retain the translates indexed by \(-M,\ldots,M\).
The real-line integral differs from \(2M+1\) times the circle integral
by a bounded endpoint error, since both the endpoint regions and the
multiplicities there are uniformly bounded. Divide by \(2M+1\) and let
\(M\to\infty\).
\end{proof}

\subsection{The spillover inequality}

Sum over ordered pairs of distinct speeds and put
\[
 W(V)=\frac1{n^2}\sum_{u\ne v\in V}\min(1,2n/H(u,v)),\qquad
 Q(r)=\sum_{u\ne v\in V}Q_r(u,v).
\]
The difference between \(\int N_r(N_r-1)\) and \(Q(r)\) is linear
in \(r\). The identity
\[
 m(m-1)-2(m-1)_+=2\sum_{k=2}^{n-1}(m-k)_+\quad(0\le m\le n,\ m\in\mathbb Z)
\]
and Lemma~\ref{lem:dilation} therefore give the second-moment spillover
inequality of Hu and Jing~\cite{HuJing}. For \(s\ge1\) and \(0\le r<1/(2s)\),
\[
 Q(sr)\ge sQ(r)+2(E_1(sr)-sE_1(r)).
\]
Lemma~\ref{lem:pair} gives \(Q(r_0)\ge1-1/n-W(V)\) and
\(Q(sr_0)\le s^2(1-1/n)\). Also \(E_1(sr_0)\ge s-1\), since the
sum of arc lengths is \(s\) and their union has measure at most \(1\).
Substituting \(s=7/5\) and rearranging yields
\begin{equation}
 g(V)\ge\frac{3}{35}-\frac{W(V)}2+\frac1{5n}.\tag{C1}\label{eq:spill}
\end{equation}

\section{Reduction to separated speed sets}\label{sec:reduce}

For \(P>0\), call a set \(P\)-separated if \(H(u,v)>P\) for every pair
of distinct elements. Suppose \(W(S)\le\delta\) for all \(P\)-separated
sets \(S\), with \(W(S)\) normalized by \(|S|\) as above.
Keep \(r_0\) fixed and, for \(S\subset V\), put
\[
 E(S)=\sum_{v\in S}|B_v(r_0)|-
           \left|\bigcup_{v\in S}B_v(r_0)\right|.
\]
If \(H(u,v)\le P\) for distinct \(u,v\in S\), deleting \(u\) decreases
\(E(S)\) by at least the common-center overlap \(1/(nP)\).
Since \(E(V)=g(V)\) and \(E(S)\ge0\), the nonempty separated set
remaining at termination has \(1-z\le Pg(V)\), where \(z=|S|/n\).
For pairs within \(S\),
\(\min(1,2n/H)\le z^{-1}\min(1,2|S|/H)\). The remaining pairs
contribute at most \(1-z^2\) to \(W(V)\). Since \(W(S)=0\) for a
singleton, in all cases
\begin{equation}
 W(V)\le zW(S)+1-z^2\le\delta+2Pg(V).\tag{C2}\label{eq:delete}
\end{equation}
We will use \(P=e^{6880}\) and \(\delta=17/100\).

\section{Stabilization of weighted graphs}\label{sec:stable}

Set \(q=201/50=4(1+\sigma)\) and \(\sigma=1/200\). We use a
bipartite graph with labeled positive integer vertices, with distinct
values on each side, and with edge set \(E\). If the side sizes are
\(m_1,m_2\), let \(m=m_1m_2\le n^2\). For scales \(a,b>0\), set
\[
 w(x,y)=\min\left(1,\frac{\gcd(x,y)}{\max(ax,by)}\right),\quad
 \rho=\frac{\sum_{(x,y)\in E}w(x,y)}m,\quad
 J=\rho^q mab(m/n^2)^\sigma.
\]
For the initial graph use two copies of \(V\), all ordered pairs of
distinct values, and \(a=b=1/(2n)\). Then
\(\rho=W(V)\) and \(J=W(V)^q/4\).

Following the valuation-layer approach to GCD graphs~\cite{GreenWalker,HuJing},
allow restriction to a nonempty subrectangle. Also allow restriction
to fixed \(p\)-adic valuation layers \(i,j\) on the two sides, followed
by division of their values by \(p^i,p^j\) and replacement of the
scales by \((ap^{(i-j)_+},bp^{(j-i)_+})\).
The latter operation preserves every retained weight and multiplies
\(ab\) by \(p^{|i-j|}\). It also preserves same-side heights.
An operation that strictly increases the score decreases the sum of
all vertex values, by removing a vertex or reducing a value. Starting
from positive edge mass, repeated improvement terminates at a state
with score at least the initial score, in which no allowed operation
increases the score.

At such a state, let \(\alpha,\beta\) be the uniform vertex measures
and let \(\mu\) be the normalized edge weight. For a restriction to
\(I\times J'\), put \(v=\alpha(I)\beta(J')\). When \(v>0\), the score
ratio is \(\mu(I\times J')^qv^{1+\sigma-q}\), so stability gives
\begin{equation}
 \mu(I\times J')\le[\alpha(I)\beta(J')]^{3/4}.\tag{L04}\label{eq:rect}
\end{equation}
The inequality is immediate when \(v=0\).
For a valuation-layer restriction the ratio has the additional factor
\(p^{|i-j|}\). Writing \(\alpha_i,\beta_j\) for the valuation marginals
and \(\mu_p\) for the edge measure on valuation pairs, we obtain
\begin{equation}
 \mu_p(i,j)\le p^{-|i-j|/q}(\alpha_i\beta_j)^{3/4}.
 \tag{L05}\label{eq:valuation}
\end{equation}
Stabilizing the fourth-power score \(\rho^4mab\) separately gives
\eqref{eq:rect} and \eqref{eq:valuation} with \(4\) in place of \(q\)
in the latter exponent, which is a stronger bound.

\section{Bilinear bounds for the Green kernel}\label{sec:kernel}

\subsection{The continuous estimate}

Let \(f\) be a nonnegative step function supported on finitely many
unit cells \([i,i+1)\), \(i\in\mathbb Z\). For \(a>0\), set
\[
 U(x)=\frac1{2a}\int_{\mathbb R}e^{-a|x-y|}f(y)\,dy .
\]
Splitting the integral at \(y=x\) shows that \(U\) is continuously
differentiable and \(U''=a^2U-f\) between the jumps of \(f\).
Both \(U\) and \(U'\) decay exponentially at infinity. Integration
by parts, with cancellation of the terms at the jumps, gives
\(E:=\int(U'^2+a^2U^2)=\int fU\).
If \(f\) has positive integral, \(U\) attains a positive maximum \(M\).
Writing \(L=\int U^4\), we have
\[
 E-\frac{a^2L}{M^2}\ge 2a\int |U'|U\sqrt{1-U^2/M^2}
\]
by the nonnegativity of the square of
\(|U'|-aU\sqrt{1-U^2/M^2}\).
On each side of a maximum point, the total variation of
\(\Phi(U)\), where \(\Phi(u)=\int_0^u s\sqrt{1-s^2/M^2}\,ds\),
is at least \(\Phi(M)=M^2/3\). Thus
\[
 E\ge a^2L/M^2+4aM^2/3\ge \frac{4a^{3/2}}{\sqrt3}\sqrt L.
\]
H\"older's inequality gives \(E\le\|f\|_{4/3}L^{1/4}\), whence
\(E\le\sqrt3\,\|f\|_{4/3}^2/(4a^{3/2})\).
For the potentials of \(f\) and \(g\), the same integration by parts gives
\(\int(U_f'U_g'+a^2U_fU_g)=\int fU_g\).
The double integral below is \(2a\int fU_g\). Cauchy--Schwarz for the
mixed energy therefore gives
\[
 \iint e^{-a|x-y|}f(x)g(y)\,dx\,dy
 \le \frac{\sqrt3}{2\sqrt a}\|f\|_{4/3}\|g\|_{4/3}.
\]
The bound holds also when either function vanishes almost everywhere.

\subsection{Discretization and the diagonal term}

Take \(\lambda=e^{-a}\). For sequences \(x_i,y_j\) with finite
support, use step functions constant on the unit intervals
\([i,i+1)\), \([j,j+1)\). The integral of the exponential kernel over
an off-diagonal cell is \(A\lambda^{|i-j|}\) and over a diagonal
cell is \(D_0\), where
\[
 A=\frac{(1-\lambda)^2}{\lambda a^2},\qquad
 D_0=2\int_0^1(1-s)e^{-as}\,ds=\frac{2(a-1+\lambda)}{a^2}.
\]
In particular, \(A-D_0=2(\sinh a-a)/a^2>0\). The continuous
estimate gives, for nonnegative sequences,
\[
 \sum_{i,j}\lambda^{|i-j|}x_i y_j
 \le C(\lambda)\|x\|_{4/3}\|y\|_{4/3}
             +d(\lambda)\sum_i x_i y_i,
\]
where
\[
 C(\lambda)=\frac{\sqrt3}{2}
               \frac{\lambda a^{3/2}}{(1-\lambda)^2},
 \qquad
 d(\lambda)=1-\frac{2\lambda(a-1+\lambda)}{(1-\lambda)^2}.
\]
Since \(\sum_i x_i y_i\le\|x\|_{4/3}\|y\|_{4/3}\), the bilinear norm
is bounded by \(\mathcal N(\lambda):=C(\lambda)+d(\lambda)\).
For probability marginals \(\alpha_i,\beta_i\), the sequences
\(\alpha_i^{3/4},\beta_i^{3/4}\) each have \(\ell^{4/3}\) norm \(1\).

\section{Concentration at large primes}\label{sec:large}

\subsection{Choosing the center and bootstrapping}

Assume \eqref{eq:rect} and \eqref{eq:valuation}. Fix a prime \(p>142\)
and put \(\lambda=p^{-1/q}\). Summing \eqref{eq:valuation} gives
a lower bound of \(1\) for the bilinear kernel sum. Since \(p\ge149\)
and \(\exp(62511/12500)<149\), we may replace \(\lambda\) in this sum
by the larger value \(\lambda_0=e^{-311/250}\). For
\(D=\sum_i(\alpha_i\beta_i)^{3/4}\), the diagonal estimate of
Section~\ref{sec:kernel} and the rational bounds in
Appendix~\ref{sec:cert} yield
\[
 1\le C(\lambda_0)+d(\lambda_0)D,\qquad
 C(\lambda_0)<0.68365,\quad 0<d(\lambda_0)<0.39442 .
\]
At \(D=401/500\), the rational upper bound on the right is
\(1-629/25000000\), so \(D\ge401/500\). Let
\(m_*=\max_i(\alpha_i+\beta_i)/2\). Termwise,
\((\alpha_i\beta_i)^{3/4}\le\sqrt{m_*}\sqrt{\alpha_i\beta_i}\).
Cauchy--Schwarz gives \(\sum_i\sqrt{\alpha_i\beta_i}\le1\), hence
\(m_*\ge D^2\). Fix an index \(k\) attaining \(m_*\), and write
\(z=1-(\alpha_k+\beta_k)/2\). Then \(z\le 89199/250000<0.357\).

Return to the actual parameter \(\lambda\) while keeping \(k\) fixed.
The center cell and the rectangle off both the center row and center
column have masses at most \((1-z)^{3/2}\) and \(z^{3/2}\), respectively.
For the remainder of the row and column, use H\"older's inequality
with \(\sum_{i\ne k}\lambda^{4|i-k|}\le2\lambda^4/(1-\lambda^4)\).
A second application to the two marginals, together with
\(\alpha_k^3+\beta_k^3\le1+(1-2z)^3\), gives
\[
 z^{1/4}b(z)\le
  \frac{2\lambda}{(1-\lambda^4)^{1/4}}[1+(1-2z)^3]^{1/4},
 \qquad b(z)=\frac{1-(1-z)^{3/2}-z^{3/2}}{z}.
\]
The case \(z=0\) requires no further estimate.
For \(0<z\le1/2\), let \(s=\sqrt z+\sqrt{1-z}\). Then
\(b(z)=2(1-z)(s+2)/(s+1)^2\), so \(b(z)/[1+(1-2z)^3]^{1/4}\)
is positive and nonincreasing. Indeed, \(s\) is increasing and
\((s+2)/(s+1)^2\) is decreasing. For the other factor
\((1-z)/[1+(1-2z)^3]^{1/4}\), its fourth power is increasing in
\(t=1-2z\), as \(4+t^3-3t^2>0\). Thus if \(z\le S\le1/2\),
\[
 z\le \frac{16\lambda^4[1+(1-2S)^3]}{(1-\lambda^4)b(S)^4}.
\]
Using \(\lambda\le144004613283/500000000000\) and the endpoints in
Table~\ref{tab:bootstrap} gives the successive upper bounds
\[
 0.357,\quad 0.342,\quad 0.312,\quad 0.262,\quad
 0.201,\quad 0.150,\quad 0.121.
\]

\subsection{The exceptional edge mass at one prime}

Let \(E_p\) be the mass of edges whose two valuations are both
strictly below \(k\) or both strictly above \(k\). Let \(z_-,z_+\)
be the averages, over the two marginals, of the masses strictly
below and strictly above \(k\), so \(z=z_-+z_+\). For \(z>0\) put
\(u=\min(z_-,z_+)/z\), and define
\[
 \begin{aligned}
 f(u)&=u^{3/4}+(1-u)^{3/4},&
 h_0(u)&=u^{3/2}+(1-u)^{3/2},\\
 j(u)&=[u(1-u)]^{3/4},&
 A_1&=(1-\lambda^4)^{-1/4}.
 \end{aligned}
\]
By \eqref{eq:rect} and the arithmetic--geometric mean inequality,
\(E_p\le z^{3/2}h_0(u)\). For tails on opposite sides of \(k\),
the kernel factors into two geometric tails, each at distance at
least one from \(k\). H\"older bounds their combined mass by
\(A_1^2\lambda^2[(\alpha_-\beta_+)^{3/4}+(\alpha_+\beta_-)^{3/4}]\),
where the subscripts denote the corresponding tails. This is at most
\(4^{3/4}A_1^2\lambda^2(z_-z_+)^{3/4}\), by Cauchy--Schwarz and
\(\sum t_i^{3/2}\le(\sum t_i)^{3/2}\) for nonnegative \(t_i\).
The four center-to-tail rectangles similarly have total mass at most
\(2^{3/4}A_1\lambda z^{3/4} f(u)[1+(1-2z)^3]^{1/4}\);
here the last factor bounds \((\alpha_k^3+\beta_k^3)^{1/4}\).
Adding the center cell accounts for all mass.

Write
\[
 A(z)=\frac{1-(1-z)^{3/2}}z,\qquad C_*(z)=[1+(1-2z)^3]^{1/4},
 \qquad h=h_0(u)+4^{3/4}A_1^2\lambda^2 j(u).
\]
The resulting inequality is
\[
 z^{1/4}[A(z)-h\sqrt z]\le 2^{3/4}A_1\lambda f(u)C_*(z).
\]
For \(0<z\le1/8\) and \(\lambda\le1/3\), one has
\(1/\sqrt2\le h\le2\).
The identity \(A(z)=\tfrac32\int_0^1\sqrt{1-tz}\,dt\) gives
\(1\le A\le3/2\) and \(A'\le-3/8\).
Thus \(A(z)-h\sqrt z>0\), its derivative is at most \(-11/8\),
and its logarithmic derivative is at most \(-11/12\).
Also \(C_*'/C_*\ge-3/4\), by \(2t^2\le1+t^3\) for \(0\le t\le1\).
The ratio \([A(z)-h\sqrt z]/C_*(z)\) is therefore nonincreasing.
For \(z\le S\le1/8\) define
\[
 B(S,\lambda,u)=
 \frac{2^{3/4}A_1 f(u) C_*(S)}{A(S)-h\sqrt S}.
\]
We obtain
\[
 z\le\lambda^4 B(S,\lambda,u)^4,\qquad
 E_p\le p^{-100/67}B(S,\lambda,u)^6 h_0(u),
\]
using \(6/q=100/67\). If \(z=0\) then \(E_p=0\).

\subsection{Summing the exceptional masses}

Starting at \(S=0.121\), partition \([0,1/2]\) into 256 intervals.
On \([l,v]\), bound \(f(u),h_0(u),j(u)\) above by
\(f(v),h_0(l),j(v)\), respectively, using their monotonicity on
\([0,1/2]\). In each prime range, use a rational upper bound
\(\lambda_{\mathrm U}\) for \(\lambda\) in \(B\) and retain the
factor \(p^{-100/67}\) in the bound for \(E_p\).
The certificate gives a positive lower bound \(d\) on the denominator,
an upper bound \(b\) on the numerator, and verifies
\[
 b\le b_{\mathrm U}d,\qquad
 b^6 h_{0,\mathrm U}\le e_{\mathrm U}d^6,\qquad
 \lambda_{\mathrm U}^4 b_{\mathrm U}^4\le S_{\mathrm{next}}.
\]
Here \(h_{0,\mathrm U}\) bounds \(h_0\) on the cell. The rational
enclosures of radicals are checked by integer powers. These tests
give the next upper bound for \(z\) and the coefficient \(e_{\mathrm U}\)
for the exceptional mass throughout the cell.

The certified sums over \(142<p\le10^4\), \(10^4<p\le10^6\), and
\(p>10^6\) are less than \(21/50\), \(23/1000\), and \(1/125\),
respectively. The first two ranges use finite sieve tables in which
each discarded integer is accompanied by a proper divisor. On the
first range the powers are bounded individually; on the second,
subinterval counts are multiplied by an upper bound for the summand
on each subinterval. For the last range, primes are restricted to the
eight reduced classes modulo \(30\), and each progression is bounded
by its first term and an integral:
\[
 \sum_{j\ge0}(b+30j)^{-s}\le b^{-s}+\frac{b^{1-s}}{30(s-1)},
 \qquad s=100/67.
\]
The rational data for these checks are given in~\cite{LeanCode}.

The total is less than \(0.451<23/50\). Only finitely many primes
occur among the vertex values; outside this set choose the center
zero. By the union bound, a proportion at least \(r_*:=27/50\) of
the edge weight has the center between its endpoint valuations at
every prime \(p>142\). For a graph of \(142\)-rough integers, meaning
integers with no prime divisor at most \(142\), use center zero
also at all primes \(p\le142\).

\section{Removing small prime factors}\label{sec:small}

\subsection{A cut estimate and a median inequality}\label{sec:cut}

Regard the two labeled sides as disjoint, and let
\(\omega=(\mu+\mu^{\mathsf T})/2\), with marginal \(\pi\).
For a vertex subset \(S\), set \(z=(\alpha(S)+\beta(S))/2\), extending
each marginal by zero to the other side, and write
\(e=\mu(S,S^c)+\mu(S^c,S)\). By \eqref{eq:rect},
\[
 e\ge1-z^{3/2}-(1-z)^{3/2},\qquad \pi(S)\le z^{3/2}+e/2 .
\]
For \(0<z\le1/2\), the first right-hand side is at least
\((2^{3/2}-2)z^{3/2}\), by the monotonic decrease of
\([1-(1-z)^{3/2}]/z^{3/2}\) and its value at \(1/2\).
The same bound holds at \(z=0\). It follows that
\(e\ge(2-\sqrt2)\pi(S)\); for \(z\ge1/2\), use \(S^c\). Hence
\[
 \omega(S,S^c)\ge(1-1/\sqrt2)\min\{\pi(S),\pi(S^c)\}.
\]
For any nonnegative integer-valued vertex function \(F\) and a
\(\pi\)-median \(m_F\),
integrate over the cuts \(S=\{F>t\}\). The integral of the smaller
side mass is \(\mathbb E_\pi|F-m_F|\). The integral of the directed
cut mass is \(\tfrac12\mathbb E_\mu|F(x)-F(y)|\), by symmetry
of \(\omega\). Thus
\[
 \mathbb E_\pi|F-m_F|
 \le (1+1/\sqrt2)\mathbb E_\mu|F(x)-F(y)|.
\]

\subsection{A weighted valuation moment}

For each prime \(p\le142\), let \(k_p\) be a median of \(\nu_p\) and put
\(B(x)=\sum_{p\le142}|\nu_p(x)-k_p|\log p\).
For \(0<\tau_p<1/q\), \eqref{eq:valuation} and
Section~\ref{sec:kernel} give
\[
 \mathbb E_\mu e^{\tau_p|\nu_p(x)-\nu_p(y)|\log p}
 \le \mathcal N(p^{-1/q+\tau_p}).
\]
By Jensen's inequality the expected edge distance
\(\mathbb E_\mu|\nu_p(x)-\nu_p(y)|\log p\) is at most
\(\log\mathcal N(p^{-1/q+\tau_p})/\tau_p\).
The 34 choices in Table~\ref{tab:small} give a sum less than
\(99.984727\). Multiplying by the factor in
Section~\ref{sec:cut} gives \(\mathbb E_\pi B<\overline D:=170.685\).

Retain vertices with \(B(x)\le1375\). Markov's inequality bounds
the discarded \(\pi\)-mass by \(\overline D/1375\). The weight of
edges retained at both ends is thus at least the fraction
\[
 \bar t=1-2\overline D/1375=\frac{103363}{137500}.
\]
Both sides remain nonempty. If the area changes from \(m\) to \(m'\),
the score is multiplied by at least
\(\bar t^q(m'/m)^{1+\sigma-q}\ge\bar t^q\).
For the fourth-power score the corresponding factor is at least
\(\bar t^4\).

\subsection{Stripping and restabilization}

Write \(x=s_xx_{\mathrm r}\), where all prime divisors of \(s_x\)
are at most \(142\) and \(x_{\mathrm r}\) is \(142\)-rough.
With \(d_s=\gcd(s_x,s_y)\) one has
\[
 \gcd(x,y)=d_s\gcd(x_{\mathrm r},y_{\mathrm r}),\qquad
 \max(ax,by)\ge d_s\max(ax_{\mathrm r},by_{\mathrm r}).
\]
Thus replacing each vertex by its rough part at fixed scales does
not decrease the weights. For two vertices on the same side,
comparison of the reduced numerators in \(H\) gives
\[
 H(x,y)\le H(x_{\mathrm r},y_{\mathrm r})
   \exp\left(\sum_{p\le142}|\nu_p(x)-\nu_p(y)|\log p\right).
\]
For retained vertices the exponent is at most
\(B(x)+B(y)\le2750\). If the original same-side heights exceed
\(e^{6880}\), the rough heights exceed \(e^{4130}\). This also gives
injectivity of stripping on each side. The rough graph can be
stabilized again by Section~\ref{sec:stable}, preserving roughness,
separation, and the lower bound for its score.

\section{Counting rough rationals and bounding a block}

\subsection{Rational height}\label{sec:rational}

For a rough graph, use the centers from Section~\ref{sec:large} and
put \(N=\prod_p p^{k_p}\). For a vertex \(x\) write \(x/N=a_x/b_x\)
in lowest terms, and set \(x'=a_xb_x\). On a retained edge \((x,y)\),
\(k_p\) lies between \(\nu_p(x)\) and \(\nu_p(y)\) for every prime.
Consequently
\[
 x'y'=\frac{xy}{\gcd(x,y)^2},
\]
since the exponent on the left is
\(|\nu_p(x)-k_p|+|\nu_p(y)-k_p|=|\nu_p(x)-\nu_p(y)|\).
Define the height of a pair of fractions \(a/b,c/d\) to be \(H(ad,bc)\).
Dividing both arguments by \(N\) leaves their relative height unchanged.

\begin{lemma}\label{lem:rc}
Let \(R>0\), \(\ell\ge4096\), and let \(S\subset[R,2R)\) be a finite
\(e^\ell\)-separated set of positive rationals with \(142\)-rough
numerators and denominators in lowest terms. For every \(T>0\),
\[
 \#\{a/b\in S:ab\le T\}\le 1+\frac{17}{25}\frac T\ell,
\]
where the fractions in the count are in lowest terms.
\end{lemma}

\subsection{Counting numerators at a fixed denominator}

We first reduce Lemma~\ref{lem:rc} to a pair of sieve estimates.
Put \(P=e^\ell\). If \(\sqrt2 T\le P\), any two fractions in the
count have \(ad,bc<\sqrt2 T\), since \((ad)(bc)\le T^2\) and
\(1/2<ad/(bc)<2\). Separation thus allows at most one.
We may assume \(R\ge1/\sqrt2\). Indeed, if \(S\ne\varnothing\) and
\(R<1/\sqrt2\), take reciprocals and use the least reciprocal as the
new lower endpoint \(R'\). Then \(R'>1/(2R)>1/\sqrt2\), and all
reciprocals lie in \([R',2R')\). Roughness, the numerator-denominator
product, and the relative height are invariant under reciprocation.

Fix a denominator \(b\le B:=\sqrt{T/R}\). The numerators \(x\) lie in
\([Rb,U_b]\), where \(U_b=\min(2Rb,T/b)\). They are \(P\)-separated.
If \(U_b\le P\), there is at most one. More generally take positive
integers \(X,Z\) with \(Z^2\le X\) and \(XZ\le P\), and factor each
numerator as \(x=sr\), with \(s\) supported on primes at most \(X\)
and \(r\) \(X\)-rough, meaning that it has no prime divisor at most \(X\).

If \(r>U_b/X\), then \(s<X\); the same \(r\) cannot occur for two
numerators, since their height would be less than \(X\).
When \(s=1\) we have \(r\in[Rb,U_b]\), and when \(s>1\) we have
\(s\ge143\) and \(r\le U_b/143\).
If \(r\le U_b/X\), choose a divisor \(d\mid x\) with
\(1\le d\le X\) and \((x/d)Z\le U_b\).
When \(s\le Z\), take \(d=1\), since \(sZ\le X\).
Otherwise, if \(s\) has a prime factor greater than \(Z\), take that
factor. If it does not, multiply its prime factors, with multiplicity,
until the product first exceeds \(Z\); the product is at most
\(Z^2\le X\). In both cases \(d>Z\) suffices.
Assign to \(x\) the \(Z\) integers \((x/d)j\), \(1\le j\le Z\).
A collision between the images of \(x,y\) would express \(x/y\) as
a ratio of positive integers at most \(XZ\), contradicting separation.
Hence the number of numerators at denominator \(b\) is at most
\[
 N_X[Rb,U_b]+N_X[1,U_b/143]+U_b/Z,
\]
where \(N_X[I]\) counts the positive \(X\)-rough integers in \(I\).
An empty interval contributes zero.

\subsection{A finite Selberg sieve}

For an integer \(X\ge2\) and \(0<U\le A\), we use
\[
 N_X[U,A]\le\frac{A-U}{\log(X+1)}+\frac{16X^2}{\log(X+1)^2}.
\]
To see this, put
\(G=\sum_{r\le X,\ r\ {\rm squarefree}}1/\varphi(r)\).
With \(\mu_{\mathrm M}\) denoting the M\"obius function, set
\[
 w_d=\frac dG\sum_{\substack{r\le X,\ d\mid r\\ r\ {\rm squarefree}}}
         \frac{\mu_{\mathrm M}(r/d)\mu_{\mathrm M}(r)}{\varphi(r)}.
\]
Then \(w_1=1\), and the square of \(\sum_{d\mid m}w_d\) equals \(1\)
for \(X\)-rough \(m\). Sum the square over the entire interval.
The count of multiples of any least common multiple differs from
\((A-U)/[d,e]\) by at most \(1\). The main quadratic form is
\[
 \sum_{d,e}\frac{w_dw_e}{[d,e]}
 =\sum_r\varphi(r)\left(\sum_{r\mid d}\frac{w_d}{d}\right)^2
 =1/G.
\]
The first identity uses \(\sum_{r\mid(d,e)}\varphi(r)=(d,e)\).
Finite M\"obius inversion evaluates the inner sum as
\(\mu_{\mathrm M}(r)/(G\varphi(r))\) for squarefree \(r\le X\) and
as zero otherwise.

Grouping \(\sum_{m\le X}1/m\) by the squarefree kernel
\(\prod_{p\mid m}p\) shows \(G\ge\log(X+1)\): the contribution of
kernel \(r\) is at most \(\prod_{p\mid r}(p-1)^{-1}=1/\varphi(r)\).
For the error, the terms in the formula for \(w_d\) have a common
sign for fixed \(d\), so
\[
 \sum_d|w_d|=\frac1G\sum_{\substack{r\le X\\r\ {\rm squarefree}}}
        \prod_{p\mid r}\left(1+\frac2{p-1}\right)
 \le \frac XG\prod_{p\le X}\left(1+\frac2{p(p-1)}\right)
 \le \frac{4X}G.
\]
For the first inequality, expand the product after removing the
squarefree restriction and bound the number of multiples of \(d\)
by \(X/d\). For the last inequality, the factor at \(2\) is \(2\);
for odd primes \(p=2k+1\), use
\(2/(p(p-1))\le 2/[3k(k+1)]\) and \(e^{2/3}<2\).
The error is at most \((\sum_d|w_d|)^2\), as required.

\subsection{Sieve parameters and summation over denominators}

Use \(X=\lfloor e^{\alpha\ell}\rfloor\) and
\(Z=\lfloor e^{\ell/10}\rfloor\). The choices
\(\alpha=1/5,4/5,9/10\) satisfy the required packing conditions,
and \(1/Z\le2e^{-\ell/10}\), \(\log(X+1)>\alpha\ell\).

\paragraph{The case \(B\ge128\).}
Use \(\alpha=1/5\). Write \(U(x)=\min(2Rx,T/x)\) and
\(G_1(x)=U(x)-Rx\). Both functions are unimodal on \((0,B]\),
with maximum at \(B/\sqrt2\). Rough denominators lie in the
classes \(1,5\bmod6\). On a progression of mesh \(6\), the
sum of either function is at most its integral divided by \(6\)
plus its maximum. Indeed, remove the largest sampled term and
compare each other term to the integral over the mesh interval
toward that term. The two classes contribute an integral divided
by \(3\) plus two maxima. The relevant integrals and maxima are
\[
 \int_0^B U=\frac{T(1+\log2)}2,\qquad
 \int_0^B G_1=\frac{T\log2}2,\qquad
 \max U=\frac{\sqrt2 T}{B},\qquad
 \max G_1=\frac T{\sqrt2 B}.
\]
The same comparison on the integer grid gives
\(\sum_b U_b\le\int_0^B U+\max U<3T/2\).
The sieve error per denominator is at most
\(32e^{2\ell/5}/(\ell/5)^2\). Since there are at most \(B\)
denominators and \(B/T\le\sqrt2e^{-\ell/2}\), the normalized
sieve error, after multiplying by \(\ell/T\), is at most
\((1200/\ell)e^{-\ell/10}\). The normalized packing error is at
most \(3\ell e^{-\ell/10}\).
The normalized main term is at most
\[
 5\left[\frac{\log2}{6}+\frac2{\sqrt2 B}
    +\frac1{143}\left(\frac{1+\log2}{6}+\frac{2\sqrt2}{B}\right)\right].
\]
Using \(B\ge128\), \(\log2<0.7\), \(1/\sqrt2\le3/4\),
\(\sqrt2\le3/2\), and \(\ell\ge4096\), the sum of the three
bounds is less than \(0.673<17/25\).

\paragraph{The case \(B<128\).}
Only \(1\) can occur as a rough denominator.
Write \(A=\min(2R,T)\). If \(A\le e^\ell\), the count is at most
one. If \(A>e^{1.7\ell}\), use \(\alpha=4/5\) and \(A-R\le A/2\).
The sieve and packing bounds give at most
\[
 \frac A\ell\left[\frac{1/2+1/143}{4/5}
       +(50/\ell+2\ell)e^{-\ell/10}\right]<0.635\,A/\ell.
\]
If \(e^\ell<A\le e^{1.7\ell}\), choose \(\alpha=9/10\), so
\(A<X^2\) and each \(X\)-rough integer at most \(A\) is \(1\) or
prime. For \(\vartheta(x)=\sum_{p\le x}\log p\) we use
\begin{equation}
 0.92x\le\vartheta(x)\le1.107x\qquad (\log x\ge1000).
 \label{eq:chebyshev}
\end{equation}
The first rough interval has count at most
\((\vartheta(A)-\vartheta(R-1))/(\ell-1)\), using \(R-1\) to
retain the left endpoint. The second has count at most
\(1+\vartheta(A/143)/(0.9\ell-1)\).
With \(R\ge A/2\), the combined count including the packing term is
less than
\[
 1+\frac{0.647 A+0.92}{\ell-1}
       +\frac{1.107 A}{143(0.9\ell-1)}
       +2Ae^{-\ell/10}
 <1+0.657\,A/\ell.
\]
Since \(A\le T\), all cases give Lemma~\ref{lem:rc}.
\hfill\(\square\)

\paragraph{Proof of \eqref{eq:chebyshev}.}
Write \(F(x)=\log(\lfloor x\rfloor!)\) and
\(J(x)=F(x)-F(x/2)-F(x/3)-F(x/5)+F(x/30)\).
The prime-power expansion of the factorial gives the coefficient
\(j(\lfloor x/p^k\rfloor)\) of the contribution from \(p^k\), where
\(j(n)=n-\lfloor n/2\rfloor-\lfloor n/3\rfloor-\lfloor n/5\rfloor+
\lfloor n/30\rfloor\). Checking residues modulo \(30\) gives
\(j\in\{0,1\}\), with \(j(n)=1\) for \(1\le n<6\). Thus
\(J(x)\le\psi(x)\le J(x)+\psi(x/6)\), where \(\psi\) sums over prime
powers. Integral comparison for \(F\) gives
\[
 |J(x)-C_0x|\le5+10\log x\quad(x\ge30),\qquad
 C_0=\frac{\log2}{2}+\frac{\log3}{3}+\frac{\log5}{5}-\frac{\log30}{30}.
\]
For instance, \(|F(x)-(x\log x-x)|\le1+2\log x\) for \(x\ge1\),
by comparison with the integral of \(\log\) and replacement of
\(\lfloor x\rfloor\) by \(x\).
The upper recurrence gives
\(\psi(x)\le(6/5)C_0x+200(1+\log(x+1))^2\)
by strong induction on \(\lfloor x\rfloor\). For \(x<30\) the
error term alone suffices. For \(x\ge30\), the difference of the
error terms at \(x\) and \(x/6\) exceeds \(5+10\log x\).
The prime-power contribution \(\psi(x)-\vartheta(x)\) is at most
\(\sqrt{x}\log x\), since each prime at most \(\sqrt{x}\) contributes
at most \(\log x\). These estimates at \(\log x\ge1000\) and rational
enclosures for \(C_0\) give \eqref{eq:chebyshev}.

\subsection{A truncated double sum}

Let \(x'_1,\ldots,x'_a\) and \(y'_1,\ldots,y'_b\) be nonempty
lists of positive numbers with counting functions at most \(1+\kappa T\) for
every \(T>0\). Then for \(c>0\),
\[
 \sum_{i,j}\min(1,c/\sqrt{x_i'y_j'})
 \le a+b+ab\,\Psi\left(\frac{c\kappa}{\sqrt{ab}}\right),
 \qquad
 \Psi(t)=
 \begin{cases}
 4t-3t^2+2t^2\log t,&0<t\le1,\\
 1,&t\ge1.
 \end{cases}
\]
Order each list and remove its first entry. Each remaining \(i\)-th
entry is at least \(i/\kappa\). The removed row and column cost at
most \(a+b\); each remaining summand is bounded by the integral of
the corresponding decreasing function on its unit cell.
Rescaling the rectangle gives
\(\Psi(t)=\int_{[0,1]^2}\min(1,t/\sqrt{uv})\,du\,dv\).
For \(t\le1\), integrating first in \(v\) and splitting at \(u=t^2\)
gives contributions \(t^2\) and \(4t(1-t)+2t^2\log t\).
If both counting functions are at most \(\kappa T\), the same
argument needs no deletion and the term \(a+b\) may be omitted.
The lists may contain repeated values, counted with their multiplicities.

Set \(\varepsilon=1/1000\), \(t_0=19/100\), and \(r_*=27/50\).
For \(t>0\) and \(0\le\rho\le1\), the condition
\(r_*(1-\varepsilon)\rho\le\Psi(t)\) implies
\(\rho^4/t^2\le1/t_0^2\). For \(t\ge t_0\) this follows from
\(\rho\le1\). For \(0<t\le t_0\), we have
\(\log t\le-\log5<-8/5\), so \(\Psi(t)\le F(t)\), where
\(F(t)=4t-(31/5)t^2\).
Writing \(u=\sqrt{t/t_0}\in(0,1]\), we obtain
\[
 uF(t_0)-F(t)
 =u(1-u)\left(4t_0-\frac{31}{5}t_0^2(1+u+u^2)\right)\ge0 .
\]
Here \(1+u+u^2\le3\). Since \(F(t_0)=0.53618<r_*(1-\varepsilon)\),
we get \(\rho\le u\).

\subsection{A block estimate at fixed cutoff}

Consider a graph whose sides lie in \([A,2A)\) and \([B,2B)\),
each consisting of distinct \(142\)-rough integers, with side sizes
\(a,b\le n\). At cutoff \(D>0\), give an edge weight
\(\min(1,\gcd(x,y)/D)\), and let \(\rho\) be the sum of edge weights
divided by \(ab\). Put \(Q_0=AB/D^2\) and
\[
 J_{\mathrm{ph}}=\rho^q\frac{ab}{Q_0}(ab/n^2)^\sigma .
\]
Upon restricting to valuation layers \(i,j\) at \(p\) and dividing
by \(p^i,p^j\), replace \(A,B,D\) by
\(A/p^i,B/p^j,D/p^{\min(i,j)}\). The retained weights are unchanged
and \(Q_0\) is multiplied by \(p^{-|i-j|}\). Thus the stabilization
argument applies to \(J_{\mathrm{ph}}\) as well as, separately, to
\(\rho^4 ab/Q_0\), yielding \eqref{eq:rect} and \eqref{eq:valuation}.

In a stable graph of positive weight, retain the edges specified in
Section~\ref{sec:large}, of mass at least \(r_*\rho ab\).
For their common normalization \(N\), the identity of
Section~\ref{sec:rational} and \(xy<4AB\) bound the retained
weight by \(\min(1,2\sqrt{Q_0}/\sqrt{x'y'})\). Summing this upper
bound over the entire labeled rectangle is permitted.

\paragraph{A fourth-power bound without separation.}
For reduced positive fractions in \([R,2R)\) with
numerator-denominator product at most \(T\), the numerator is at
most \(\sqrt{2RT}\) and the denominator at most \(\sqrt{T/R}\).
There are at most \(\sqrt2 T\) such integer pairs.
In the stable graph for the fourth-power score, apply the
double-sum bound without deletion at \(\kappa=\sqrt2\). It gives
\(r_*\rho\le\Psi(2\sqrt{2Q_0/(ab)})\) and hence
\(\rho^4 ab/Q_0\le8/t_0^2\).
Since stabilization does not decrease the score, this holds
for every original block.

\paragraph{The \(q\)-power score with separation.}
Assume both sides are \(e^\ell\)-separated and \(\ell\ge4096\).
In the stable graph for \(J_{\mathrm{ph}}\) use
Lemma~\ref{lem:rc} with \(\kappa=(17/25)/\ell\).
If \(\min(\rho a,\rho b)\ge2/(r_*\varepsilon)\), then
\(a+b\le r_*\varepsilon\rho ab\). Subtracting this term from
the double-sum bound yields
\(r_*(1-\varepsilon)\rho\le\Psi(2\kappa\sqrt{Q_0/(ab)})\).
Thus \(\rho^4 ab/Q_0\le4\kappa^2/t_0^2\), which also bounds
\(J_{\mathrm{ph}}\).
If \(\rho a<2/(r_*\varepsilon)\), then
\(\rho^4ab/n^2\le2/(r_*\varepsilon n)\) by \(b\le n\)
and \(\rho^3\le1\), and the case of small \(\rho b\) is symmetric.
As \(q=4(1+\sigma)\), the score equals the fourth-power score
times \((\rho^4ab/n^2)^\sigma\). Use the fourth-power bound in
this case. Combining the two bounds, and including zero weight,
gives for the original block
\begin{equation}
 J_{\mathrm{ph}}\le C_{\mathrm{ph}}(\ell,n):=
 \frac4{t_0^2}\left(\frac{17}{25\ell}\right)^2
 +\frac8{t_0^2}\left(\frac2{r_*\varepsilon n}\right)^\sigma .
 \label{eq:block}
\end{equation}

\section{Global estimate and proof of the theorem}

\subsection{Dyadic decomposition}

In the rough graph of Section~\ref{sec:small}, partition \(ax\)
and \(by\) using the common grid \([h2^i,h2^{i+1})\), \(i\ge0\),
with \(h=\min(a,b)\). Both scaled speeds are at least \(h\).
Put \(A_i=h2^i/a\), \(B_j=h2^j/b\), and
\(D_{ij}=\max(aA_i,bB_j)\). On this block the original weight
is bounded by the fixed-cutoff weight at \(D_{ij}\), and
\[
 \frac{A_i B_j}{D_{ij}^2}
 =\frac1{ab}\exp\bigl(-|\log(aA_i)-\log(bB_j)|\bigr).
\]
Write \(m_i,n_j\) for the respective block sizes and \(m_L,m_R\)
for the full side sizes. Taking a \(q\)-th root in \eqref{eq:block}
bounds the block's total weight by
\[
 C_{\mathrm{ph}}(\ell,n)^{1/q}(ab)^{-1/q}n^{2\sigma/q}
 (m_i n_j)^{3/4}(2^{-1/q})^{|i-j|},
\]
because \(1-(1+\sigma)/q=3/4\).
The sequences \((m_i/m_L)^{3/4}\) and \((n_j/m_R)^{3/4}\)
have unit \(\ell^{4/3}\) norms. Applying Section~\ref{sec:kernel}
to the sum, and returning to the definition of \(J\), gives
\[
 J\le \overline B^{\,q} C_{\mathrm{ph}}(\ell,n),
 \qquad \mathcal N(2^{-1/q})\le\overline B:=2.137859 .
\]

\subsection{The bound for separated sets}

Suppose \(V\) is \(e^{6880}\)-separated. If \(n<e^{6880}/40\),
then every ordered pair contributes a summand less than \(1/20\)
in \(W(V)\), so \(W(V)<1/20\). Otherwise, the initial score, the
retention of at least the fraction \(\bar t^q\) of the score,
and the resulting separation parameter \(\ell=4130\) give
\[
 \frac{\bar t^q W(V)^q}{4}
   \le\overline B^{\,q} C_{\mathrm{ph}}(4130,n).
\]
Since \(C_{\mathrm{ph}}(4130,n)\) is decreasing in \(n\), use
\(n\ge e^{6880}/40\). The certificates give
\[
 C_{\mathrm{ph}}(4130,e^{6880}/40)<\overline C
       :=\frac{3004}{10^9},
 \qquad
 4\overline C\,\overline B^{\,q}<(0.1699\bar t)^q .
\]
Raising the second comparison to the fiftieth power reduces it
to a rational inequality. It follows that \(W(V)<0.1699<17/100\).

\subsection{Increasing the forbidden radius}

\begin{proof}[Proof of Theorem~\ref{thm:main}]
Apply \eqref{eq:delete} with \(P=e^{6880}\) and \(\delta=17/100\),
and combine it with \eqref{eq:spill} to obtain
\[
 (1+e^{6880})g(V)\ge \frac3{35}-\frac{17}{200}=\frac1{1400}.
\]
Since \(\sum_{j=0}^{20}8^j/j!>2800\), we have \(e^8>2800\) and
\(1400(1+e^{6880})<e^{6888}\). Thus \(g(V)>e^{-6888}\).

Set \(c=e^{-6888}\) and increase the radius from \(r_0\) to
\((1+c)/(2n)\). For each speed the forbidden measure increases
by \(c/n\), so the total increase in the measure of the union
is at most \(c\). Its complement still has positive measure.
Any point in the complement, represented in \([0,1)\), lies outside
every closed forbidden set and satisfies the strict inequalities
of the theorem.
\end{proof}

\appendix
\section{Numerical certificates and formalization}\label{sec:cert}

\subsection{Rational data}

For each row of Table~\ref{tab:bootstrap}, the rational numbers \(u,v\)
are upper bounds on \(\sqrt S,(1-S)^{3/2}\), respectively, so that
\((1-v)/S-u\) is a lower bound on \(b(S)\). The output \(T\) is an
upper bound on the next value of \(z\).

\begin{table}[htbp]
\centering
\small
\caption{Rational endpoint certificates for the six iterations.}
\label{tab:bootstrap}
\begin{tabular}{@{}rrrr@{}}
\toprule
Input \(S\) & Output \(T\) & Upper bound \(u\) & Upper bound \(v\)\\
\midrule
357/1000 & 171/500 & 5974947699/10000000000 & 1031208431/2000000000 \\
171/500 & 39/125 & 5848076607/10000000000 & 1067502341/2000000000 \\
39/125 & 131/500 & 2792848009/5000000000 & 5706668661/10000000000 \\
131/500 & 201/1000 & 5118593557/10000000000 & 6339931167/10000000000 \\
201/1000 & 3/20 & 896660471/2000000000 & 3571002657/5000000000 \\
3/20 & 121/1000 & 3872983347/10000000000 & 7836612789/10000000000 \\
\bottomrule
\end{tabular}
\end{table}

Table~\ref{tab:small} gives the tilting parameters and certified
upper bounds on \(\log\mathcal N(p^{-1/q+\tau_p})/\tau_p\).

\begingroup\small
\begin{longtable}{@{}rrr@{}}
\caption{Tilts and weighted valuation distances at small primes.}
\label{tab:small}\\
\toprule
Prime \(p\) & Tilt \(\tau_p\) & Rational upper bound\\
\midrule
\endfirsthead
\toprule
Prime \(p\) & Tilt \(\tau_p\) & Rational upper bound\\
\midrule
\endhead
\bottomrule
\endfoot
2 & 37 / 201 & 7670990696489 / 1000000000000 \\
3 & 347 / 2010 & 640388639331 / 100000000000 \\
5 & 107 / 670 & 662523832733 / 125000000000 \\
7 & 61 / 402 & 2365565222337 / 500000000000 \\
11 & 142 / 1005 & 4090203987259 / 1000000000000 \\
13 & 277 / 2010 & 155231431817 / 40000000000 \\
17 & 133 / 1005 & 3569839071111 / 1000000000000 \\
19 & 131 / 1005 & 1724576903557 / 500000000000 \\
23 & 17 / 134 & 3252068305893 / 1000000000000 \\
29 & 41 / 335 & 151444678617 / 50000000000 \\
31 & 122 / 1005 & 2967702979251 / 1000000000000 \\
37 & 119 / 1005 & 2811502620707 / 1000000000000 \\
41 & 39 / 335 & 681198415611 / 250000000000 \\
43 & 233 / 2010 & 671373512687 / 250000000000 \\
47 & 23 / 201 & 1306830035577 / 500000000000 \\
53 & 113 / 1005 & 314968774063 / 125000000000 \\
59 & 223 / 2010 & 2438837267921 / 1000000000000 \\
61 & 37 / 335 & 603558459043 / 250000000000 \\
67 & 22 / 201 & 117317759673 / 50000000000 \\
71 & 109 / 1005 & 230537391869 / 100000000000 \\
73 & 217 / 2010 & 285750766663 / 125000000000 \\
79 & 43 / 402 & 223184591259 / 100000000000 \\
83 & 107 / 1005 & 2198656928103 / 1000000000000 \\
89 & 106 / 1005 & 2152619194079 / 1000000000000 \\
97 & 7 / 67 & 1048611830883 / 500000000000 \\
101 & 209 / 2010 & 258966201167 / 125000000000 \\
103 & 104 / 1005 & 1029736157291 / 500000000000 \\
107 & 69 / 670 & 2035878906713 / 1000000000000 \\
109 & 69 / 670 & 2024512845143 / 1000000000000 \\
113 & 103 / 1005 & 2002582458967 / 1000000000000 \\
127 & 203 / 2010 & 483303801639 / 250000000000 \\
131 & 101 / 1005 & 957616588461 / 500000000000 \\
137 & 1 / 10 & 188956832491 / 100000000000 \\
139 & 1 / 10 & 734896491 / 390625000 \\
\end{longtable}
\endgroup

\subsection{Verification and source locations}

The rational power tests reduce to integer powers. For example,
for positive integers \(a,b,n\), the inequality
\(a/b\ge n^{-100/67}\) is equivalent to \(a^{67}n^{100}\ge b^{67}\).
Logarithms are enclosed by truncating
\[
 \log x=2\sum_{j\ge0}\frac{v^{2j+1}}{2j+1},
 \qquad v=\frac{x-1}{x+1},
\]
and bounding the remainder by a geometric series. The large-prime
interval certificates and sieve tables are supplied with the Lean
material~\cite{LeanCode}. Module names in Table~\ref{tab:lean}
refer to files under \texttt{LonelyRunnerNew/}.
The main theorem is the declaration \texttt{LonelyRunnerNew.lonely\_runner\_gap}
in \texttt{SpilloverInputProved.lean}.

\begingroup\small
\begin{longtable}{@{}>{\raggedright\arraybackslash}p{0.25\linewidth}
                       >{\raggedright\arraybackslash}p{0.70\linewidth}@{}}
\caption{Proof steps and corresponding Lean modules.}\label{tab:lean}\\
\toprule
Step & Module or declaration\\
\midrule
\endfirsthead
\toprule
Step & Module or declaration\\
\midrule
\endhead
\bottomrule
\endfoot
Covering and deletion &
\texttt{TwoSpeedIntersection}, \texttt{PeriodicLevelDilation},
\texttt{SpilloverInputProved}, \texttt{ForestInputProved}\\
Stabilization and kernel &
\texttt{InitialTaggedState}, \texttt{ValuationDistribution},
\texttt{GreenKernel}, \texttt{KernelInputProved}\\
Large-prime exceptions &
\texttt{FixedScaleNumerics}, \texttt{BootstrapChain},
\texttt{LCNumericChecker}, \texttt{LCFiniteBudgets},
\texttt{LCQuantitativeRetention}\\
Removing small factors &
\texttt{SmallPrimeRowsProved}, \texttt{CoreCut},
\texttt{JointCore}, \texttt{OriginalSeparatedPeeling}\\
Sieve and block bound &
\texttt{FiniteIntervalSieve}, \texttt{ChebyshevExplicit},
\texttt{RationalCountProved}, \texttt{SaturatedTerminal}, \texttt{RoughBlockProved}\\
Global bound and theorem &
\texttt{DyadicPhysicalPartition}, \texttt{DyadicKernelCertificate},
\texttt{DyadicRoughAssembly}, \texttt{NumericEndgame},
\texttt{SpilloverInputProved}\\
\end{longtable}
\endgroup
\begin{thebibliography}{9}
\bibitem{HuJing}
X. Hu and Y. Jing,
\emph{Near optimal gap in the Lonely Runner Conjecture},
\href{https://arxiv.org/abs/2610.23790v1}{arXiv:2610.23790v1}, 2026.

\bibitem{GreenWalker}
B. Green and A. Walker,
\emph{Extremal problems for GCDs},
Combin. Probab. Comput. \textbf{30} (2021), 922--929.
\href{https://arxiv.org/abs/2012.02078}{arXiv:2012.02078}.

\bibitem{LeanCode}
\emph{LonelyRunner}, Lean formalization and rational certificates,
\texttt{uniform-gap/}, source snapshot at commit
\href{https://github.com/Orc-Lee-Chem/LonelyRunner/tree/2d0f3ded605f4929e0313fa4edc362064c3a1bd6/uniform-gap}{\texttt{2d0f3ded605f}},
using Lean 4.34.1. The accompanying
\texttt{candidate-source-version.json} records the per-file hashes.
\end{thebibliography}
\end{document}
